% ------------------------------------------------------------------
%  Capitolo 23 — Insegnare con la scienza dell'apprendimento
%  Parte VI
%  Ricerca di riferimento: ricerca 01 §9; 02 §9.3–9.4; ricerca 08 (da fare)
%  Voci bibliografiche principali: sweller1988, kirschner2006, kalyuga2003, rosenshine2012, freeman2014, theobald2020, deslauriers2019, eef2025
%  Epigrafe: ◐ verificata indirettamente (vedi ricerca 03 e registro, ricerca 00)
%  Scheda del capitolo: ricerca/06_indice-ragionato.md, capitolo 23
% ------------------------------------------------------------------
\chapter{Insegnare con la scienza dell'apprendimento}
\label{cap:docenti}

\epigrafe[la mente non ha bisogno di essere riempita come un vaso, ma di essere accesa come legna]{οὐ γὰρ ὡς ἀγγεῖον ὁ νοῦς ἀποπληρώσεως ἀλλ᾽ ὑπεκκαύματος μόνον ὥσπερ ὕλη δεῖται}{Plutarco, Come si deve ascoltare 48c}

\iniziale{Q}{uesto capitolo} \segnaposto{apertura: da scrivere — vedi la scheda nell'indice ragionato}.

\section{Che cosa sanno i docenti universitari}

\segnaposto{da scrivere}

\section{La lezione: istruzione esplicita e carico cognitivo}

\segnaposto{da scrivere}

\section{Apprendimento attivo in aula}

\segnaposto{da scrivere}

\section{Il metodo di studio si insegna nei corsi}

\segnaposto{da scrivere}

\section{Imparare a insegnare: il faculty development}

\segnaposto{da scrivere}

\begin{daricordare}
\segnaposto{il principio del capitolo in due righe}
\end{daricordare}

\begin{domande}
  \item \segnaposto{domanda di recupero su questo capitolo}
  \item \segnaposto{domanda di applicazione a una materia che studi}
  \item \segnaposto{domanda di ripresa su un capitolo precedente}
\end{domande}

\begin{sintesi}
  \item \segnaposto{punto di sintesi}
\end{sintesi}

\finecapitolo
